\documentclass[11pt]{article}

\usepackage[utf8]{inputenc}
\usepackage[T1]{fontenc}
\usepackage{amsmath,amssymb,amsfonts}
\usepackage{booktabs}
\usepackage{graphicx}
\usepackage{hyperref}
\usepackage[margin=1in]{geometry}
\usepackage{algorithm}
\usepackage{algorithmic}

\title{Iterative Self-Refinement for Multi-Camera Calibration\\via Learned Pose Conditioning}
\author{}
\date{}

\begin{document}
\maketitle

\begin{abstract}
We present a method for improving multi-camera extrinsic calibration using iterative self-conditioning with a learned geometry model. Starting from Pi3X, a permutation-equivariant feed-forward network that estimates camera poses from images, we observe that feeding its own predictions back as conditioning improves rotation estimates but degrades translations due to Sim(3) scale ambiguity. We propose a hybrid fusion strategy that decouples rotation and translation channels: rotation edges are extracted from conditioned inference while translation edges are preserved from unconditioned inference. Combined with Monte Carlo noise perturbation, this achieves 28\% rotation improvement and 21\% relative pose improvement over the baseline, with no external calibration data required.
\end{abstract}

\section{Introduction}
\label{sec:introduction}

We address the problem of extrinsic calibration for 8 static cameras deployed in a cattle feedlot. The scene presents challenges where classical calibration methods fail: textureless dynamic subjects (black cattle), minimal inter-camera overlap, and no physical calibration targets.

We use Pi3X~\cite{pi3x}, a permutation-equivariant feed-forward geometry model that processes all $N$ views jointly and outputs camera poses up to Sim(3) ambiguity. By aggregating predictions across 91 timesteps, we obtain an initial calibration estimate with $\sim$4.4° mean rotation error.

\textbf{Key insight.} Pi3X accepts optional pose conditioning---feeding known camera poses alongside images improves its predictions. We exploit this by feeding Pi3X's \emph{own output} back as conditioning, creating an iterative self-refinement loop. However, na\"ive iteration improves rotation but degrades translation due to Sim(3) scale ambiguity compounding across iterations.

\textbf{Contribution.} We propose a hybrid fusion strategy that:
\begin{enumerate}
    \item Decouples rotation and translation in the self-refinement loop
    \item Uses Monte Carlo noise perturbation to explore the posterior over poses
    \item Achieves 28\% rotation improvement with no external calibration data
\end{enumerate}

\section{Method}
\label{sec:method}

Our method operates in four stages. Stage~1 applies standard geometric post-processing to establish a baseline estimate. Stages~2--4 constitute our contribution: self-conditioned iterative refinement with decoupled rotation-translation fusion.

\subsection{Preliminaries: Geometric Post-Processing (Stage 1)}
\label{sec:stage1}

Given $N_t = 91$ timesteps of $N$-camera images, we run Pi3X to obtain per-timestep poses $\{T_k^c\}_{k=1}^{N_t}$ for each camera $c \in \{1,\ldots,N\}$, where each $T_k^c \in \text{SE}(3)$ is defined up to a per-timestep Sim(3) ambiguity (i.e., each timestep's output coordinate frame is arbitrary in position, orientation, and scale). We aggregate these into a single estimate via three standard steps.

\paragraph{Gauge fixing.} Since each timestep uses an arbitrarily chosen global frame, we align all timesteps to a reference via Sim(3) Procrustes~\cite{umeyama1991}:
\begin{equation}
    \{s^*, R^*, \mathbf{t}^*\} = \argmin_{s, R, \mathbf{t}} \sum_{c=1}^N \| s R \, \mathbf{p}_k^c + \mathbf{t} - \mathbf{p}_{\text{ref}}^c \|^2
\end{equation}
where $\mathbf{p}_k^c \in \mathbb{R}^3$ denotes the position of camera $c$ at timestep $k$, and $s \in \mathbb{R}^+$, $R \in \text{SO}(3)$, $\mathbf{t} \in \mathbb{R}^3$ parametrise the Sim(3) transform.

\paragraph{Outlier rejection.} For each camera, we project its $N_t$ estimated positions onto the dominant 2D plane (via PCA on camera centres) and reject outliers whose in-plane distance exceeds $3\times$ the Median Absolute Deviation (MAD). Inlier rotations are aggregated via the Fr\'{e}chet mean on SO(3)~\cite{moakher2002}---the rotation minimising the sum of squared geodesic distances to all inlier samples.

\paragraph{Pose graph optimization.} We extract pairwise relative poses $T_{ij} = (T^i)^{-1} T^j$ from inlier timesteps, computing the Fr\'{e}chet mean rotation and mean translation for each camera pair $(i, j)$. These relative poses serve as edges in a pose graph, weighted by the inverse variance of the per-timestep estimates (edges with higher consistency receive more weight). We then optimise globally via $L_1$ Weiszfeld rotation averaging on SO(3)~\cite{hartley2013} followed by translation averaging with hard co-location constraints for cameras sharing a physical mounting pole.

\paragraph{Residual error analysis.} This pipeline reduces ATE by 10\% (0.177 $\rightarrow$ 0.161) but leaves ARE unchanged (4.37\degree{} $\rightarrow$ 4.39\degree{}). The persistence of rotation error under robust averaging over 91 independent timesteps indicates that the residual is systematic model bias rather than sample noise. Since the scene provides weak geometric signal (approximately planar static structure, dynamic repetitive content), further classical aggregation cannot improve upon this estimate. This motivates the use of Pi3X's conditioning interface to inject prior geometric information, as developed in the following sections.

\subsection{Self-Conditioned Inference (Stage 2)}
\label{sec:stage2}

Pi3X supports multimodal conditioning: given approximate poses $\hat{T}^c$ and camera intrinsic matrices $\boldsymbol{\Pi}^c \in \mathbb{R}^{3 \times 3}$, it refines its predictions by attending to the provided geometric priors during encoding. We condition on the Stage~1 output:
\begin{equation}
    \{T_k^{c,(1)}\} = f_\theta(\text{imgs}_k \mid \hat{T}^c, \boldsymbol{\Pi}^c), \quad k = 1, \ldots, N_t
\end{equation}
where $f_\theta$ denotes Pi3X and $\hat{T}^c$ is the Stage~1 estimate converted to the model's input camera ordering.

Applying the same geometric post-processing to $\{T_k^{c,(1)}\}$ yields a refined estimate with ARE reduced from $4.39\degree$ to $3.49\degree$ (20\% improvement) but ATE increased from 0.161 to 0.166 (3\% degradation). Iterating this process na\"ively amplifies both effects: rotation improves further (to ${\sim}3.3\degree{}$) while translation degrades monotonically (to 0.195 ATE after 5 iterations). Section~\ref{sec:theory} provides a formal explanation: the rotation component contracts toward a fixed point on SO(3), while translations diverge due to Sim(3) scale mismatch in the conditioning pathway.

\subsection{Decoupled Hybrid Fusion (Stage 3)}
\label{sec:stage3}

Since conditioned inference improves rotation but degrades translation, we decouple them in the pose graph by sourcing edges from different inference runs. Let $R_{ij}$ and $\mathbf{t}_{ij}$ denote the relative rotation and translation from camera $i$ to camera $j$, computed as the mean over inlier timesteps. We construct fused edges:

\begin{equation}
    \tilde{T}_{ij}^{\text{fused}} = \begin{pmatrix} R_{ij}^{\text{cond}} & \mathbf{t}_{ij}^{\text{uncond}} \\ \mathbf{0}^\top & 1 \end{pmatrix}
\end{equation}

where $R_{ij}^{\text{cond}}$ is the mean relative rotation extracted from conditioned inference (Stage~2) and $\mathbf{t}_{ij}^{\text{uncond}}$ is the mean relative translation from the original unconditioned run (Stage~1). Edge weights $w_{ij}$ (inverse variance of the per-timestep relative pose estimates) are taken as $w_{ij} = \max(w_{ij}^{\text{cond}}, w_{ij}^{\text{uncond}})$.

The initial pose for graph optimization similarly uses conditioned rotations with unconditioned translations:
\begin{equation}
    T_{\text{init}}^c = \begin{pmatrix} R_{\text{cond}}^c & \mathbf{t}_{\text{uncond}}^c \\ \mathbf{0}^\top & 1 \end{pmatrix}
\end{equation}

This yields ARE = 3.45\degree{} (rotation improved) with ATE = 0.159 (translation preserved)---capturing the benefits of both sources.

\subsection{Iterated Hybrid Refinement}
\label{sec:stage4}

Since a single conditioned pass reduces rotation error by 20\%, a natural question is whether iterating the process yields further gains. We apply hybrid fusion iteratively: at each step, the optimised rotation estimate from the previous iteration serves as the new conditioning input.

\begin{algorithm}[t]
\caption{Iterated Hybrid Refinement}
\label{alg:itermc}
\begin{algorithmic}[1]
\REQUIRE Pipeline estimate $\hat{T}^c$, images, intrinsics $\boldsymbol{\Pi}^c$
\REQUIRE Iterations: $M$
\STATE $T_{\text{curr}} \leftarrow \hat{T}^c$
\STATE Compute unconditioned edges $\{\mathbf{t}_{ij}^{\text{uncond}}, w_{ij}^t\}$ \COMMENT{fixed throughout}
\FOR{$m = 1$ to $M$}
    \STATE $\{T_{k}^{c,(m)}\} \leftarrow f_\theta(\text{imgs} \mid T_{\text{curr}}, \boldsymbol{\Pi}^c)$ for all $k$
    \STATE Extract rotation edges $\{R_{ij}^{\text{cond}}, w_{ij}^R\}$ from pipeline-processed conditioned poses
    \STATE Fuse: $\tilde{T}_{ij} = (R_{ij}^{\text{cond}}, \mathbf{t}_{ij}^{\text{uncond}})$
    \STATE $T_{\text{opt}} \leftarrow \text{PoseGraph}(\bar{T}_{\text{init}}, \tilde{T}_{ij}, w_{ij})$
    \STATE $T_{\text{curr}} \leftarrow (R_{\text{opt}}, \mathbf{t}_{\text{uncond}})$ \COMMENT{anchor translations}
\ENDFOR
\RETURN $T_{\text{opt}}$
\end{algorithmic}
\end{algorithm}

\paragraph{Design choices.}
\begin{itemize}
    \item \textbf{Translation anchoring}: Translations are \emph{never} updated from conditioned outputs---only from the original unconditioned baseline. This is the key invariant that prevents Sim(3) drift.
    \item \textbf{Early stopping}: Terminate when ARE improvement between iterations falls below $0.1\degree$, or after $M=2$ iterations (empirically sufficient; see Section~\ref{sec:ablation}).
\end{itemize}

\paragraph{Monte Carlo extension.} The iterative procedure admits a natural stochastic extension: perturbing the conditioning poses before each forward pass and pooling the resulting observations across $S$ samples. Under a posterior sampling interpretation (Section~\ref{sec:theory}), this estimates the posterior mean via Monte Carlo integration, with variance decreasing as $O(1/\sqrt{S})$. We implement this with structured perturbation (shared global rotation preserving inter-camera geometry), Sim(3) alignment across samples, and pool-before-pipeline aggregation. However, our experiments (Section~\ref{sec:ablation_S}) demonstrate that this extension provides no benefit: the model's posterior under self-conditioning is so tightly concentrated that $S=1$ (deterministic iteration) is optimal.

\section{Results}
\label{sec:results}

\subsection{Experimental Setup}

\paragraph{Dataset.} 8 static cameras (2688$\times$1520) in a cattle feedlot, 91 timesteps sampled at interval 10 from synchronized video. Ground truth poses from COLMAP.

\paragraph{Evaluation metrics.}
\begin{itemize}
    \item \textbf{ATE}: Absolute Translation Error (RMSE after Sim(3) alignment)
    \item \textbf{ARE}: Absolute Rotation Error (mean angular error in degrees)
    \item \textbf{RPE rot/trans}: Relative Pose Error for all camera pairs
    \item \textbf{C2C3, C5C6}: Distance between co-located camera pairs (GT $\approx$ 0)
\end{itemize}

\subsection{Comprehensive Comparison}

\begin{table}[h]
\centering
\caption{Comprehensive comparison of all methods. ATE and RPE trans are in arbitrary (Sim(3)-aligned) units. Cost is relative to a single unconditioned Pi3X inference pass over 91 timesteps ($\sim$100s on an A100).}
\label{tab:comparison}
\begin{tabular}{@{}lcccccc@{}}
\toprule
Method & ATE$\downarrow$ & ARE$\downarrow$ & RPE rot$\downarrow$ & RPE trans$\downarrow$ & C2C3$\downarrow$ & Cost \\
\midrule
1. Raw Pi3X (quat average) & 0.177 & 4.37° & 6.26° & 0.395 & 0.088 & 1$\times$ \\
2. + Pipeline & 0.161 & 4.39° & 6.26° & 0.375 & 0.000 & 1$\times$ \\
3. + Iterative Refinement & 0.166 & 3.49° & 4.93° & 0.384 & 0.000 & 2$\times$ \\
4. + Hybrid Refinement & 0.159 & 3.45° & 4.93° & 0.322 & 0.000 & 2$\times$ \\
\textbf{5. + Hybrid MC (3$\times$4)} & \textbf{0.160} & \textbf{3.13°} & \textbf{4.57°} & \textbf{0.311} & \textbf{0.000} & 13$\times$ \\
\bottomrule
\end{tabular}
\end{table}

\subsection{Summary of Improvements}

Comparing the full method (row 5) against raw Pi3X (row 1):
\begin{itemize}
    \item \textbf{ARE}: 4.37° $\rightarrow$ 3.13° (\textbf{28\% improvement})
    \item \textbf{RPE rot}: 6.26° $\rightarrow$ 4.57° (\textbf{27\% improvement})
    \item \textbf{RPE trans}: 0.395 $\rightarrow$ 0.311 (\textbf{21\% improvement})
    \item \textbf{ATE}: 0.177 $\rightarrow$ 0.160 (\textbf{10\% improvement})
    \item \textbf{Co-location}: 0.088 $\rightarrow$ 0.000 (hard constraint enforced)
\end{itemize}

The hybrid refinement (row 4) offers the best cost-benefit tradeoff: 2$\times$ compute for 21\% rotation and 18\% RPE improvement.

\section{Ablation Studies}
\label{sec:ablation}

\subsection{Why Hybrid Fusion is Necessary}

Na\"ive iterative refinement (feeding full conditioned output back without decoupling) degrades translation while improving rotation:

\begin{table}[h]
\centering
\caption{Na\"ive vs.\ hybrid iterative refinement over 5 iterations.}
\label{tab:naive_vs_hybrid}
\begin{tabular}{@{}lcccc@{}}
\toprule
Method & ATE$\downarrow$ & ARE$\downarrow$ & $\Delta$ATE & $\Delta$ARE \\
\midrule
Unconditioned baseline & 0.161 & 4.39° & --- & --- \\
Na\"ive self-refine (iter 5) & 0.195 & 3.60° & +21\% & $-$18\% \\
Hybrid MC (iter 2) & 0.160 & 3.13° & $-$1\% & $-$29\% \\
\bottomrule
\end{tabular}
\end{table}

The Sim(3) scale ambiguity in Pi3X outputs means that conditioned translations accumulate bias across iterations. The hybrid approach avoids this by anchoring translations to the unconditioned source.

\subsection{Noise Robustness of Pose Conditioning}

We test Pi3X's conditioning with ground truth poses corrupted by varying levels of Gaussian noise:

\begin{table}[h]
\centering
\caption{Effect of conditioning pose quality on Pi3X output (after full pipeline). Translation noise is isotropic Gaussian in the scene coordinate system.}
\label{tab:noise_robustness}
\begin{tabular}{@{}cccc@{}}
\toprule
Rot.\ noise (°) & Trans.\ noise & ATE$\downarrow$ & ARE$\downarrow$ \\
\midrule
0 (perfect GT) & 0 & 0.097 & 1.22° \\
1 & 0.05 & 0.087 & 1.30° \\
2 & 0.10 & 0.080 & 1.60° \\
5 & 0.20 & 0.077 & 3.34° \\
10 & 0.50 & 0.368 & 180° \\
20 & 1.00 & 0.951 & 180° \\
\bottomrule
\end{tabular}
\end{table}

Key observations:
\begin{itemize}
    \item Conditioning tolerates up to $\sim$5° of rotation noise
    \item Beyond 10°, the conditioning signal becomes misleading and breaks the model
    \item Our self-conditioned prior ($\sim$4.4° error) is at the edge of the useful regime, explaining the diminishing returns beyond 2--3 iterations
\end{itemize}

\subsection{Convergence Behavior}

The iterated hybrid MC method converges within 2--3 iterations:

\begin{table}[h]
\centering
\caption{Convergence of iterated hybrid MC refinement (4 samples per iteration).}
\label{tab:convergence}
\begin{tabular}{@{}ccccc@{}}
\toprule
Iteration & ATE & ARE (°) & RPE rot (°) & RPE trans \\
\midrule
0 (unconditioned) & 0.161 & 4.39 & 6.26 & 0.375 \\
1 & 0.162 & 3.43 & 4.88 & 0.321 \\
2 & 0.160 & 3.13 & 4.58 & 0.309 \\
3 & 0.163 & 3.38 & 4.70 & 0.307 \\
4 & 0.160 & 180° & 4.78 & 8.55 \\
\bottomrule
\end{tabular}
\end{table}

Iteration 2 is the sweet spot. Beyond iteration 3, the conditioned rotations begin to drift and can cause orientation flips (iter 4). We recommend stopping after 2--3 iterations or when ARE improvement is $<0.1°$.

\subsection{Sim(3) Constraint Does Not Help}

We tested Sim(3)-aligning the conditioned output back to the conditioning frame before each iteration (to prevent scale drift). Results are identical to na\"ive iteration, confirming that translation degradation is not caused by scale drift but by systematic bias in how Pi3X uses pose conditioning:

\begin{table}[h]
\centering
\caption{Sim(3)-constrained vs.\ na\"ive iteration (5 iterations).}
\label{tab:sim3}
\begin{tabular}{@{}ccccc@{}}
\toprule
Iter & Sim3 ATE & Sim3 ARE & Na\"ive ATE & Na\"ive ARE \\
\midrule
0 & 0.152 & 4.69° & 0.152 & 4.69° \\
1 & 0.160 & 3.50° & 0.160 & 3.50° \\
2 & 0.180 & 3.31° & 0.179 & 3.29° \\
5 & 0.197 & 3.61° & 0.195 & 3.59° \\
\bottomrule
\end{tabular}
\end{table}

\section{Theoretical Interpretation}
\label{sec:theory}

\subsection{Contraction Mapping on SO(3)}

Let $f: \text{SO}(3)^N \rightarrow \text{SO}(3)^N$ denote the mapping from conditioning rotations to Pi3X's output rotations (after pipeline processing). If $f$ is a contraction, i.e.,
\begin{equation}
    d_{\text{SO}(3)}(f(R), R^*) \leq \lambda \cdot d_{\text{SO}(3)}(R, R^*)
\end{equation}
for some $\lambda < 1$ and true rotations $R^*$, then by the Banach fixed-point theorem, iteration $R^{(n+1)} = f(R^{(n)})$ converges to a fixed point.

\paragraph{Empirical evidence.} From the noise robustness experiment:
\begin{itemize}
    \item Input error 4.4° $\rightarrow$ output error 3.5° ($\lambda \approx 0.80$)
    \item Input error 5.0° $\rightarrow$ output error 3.3° ($\lambda \approx 0.67$)
    \item Input error 2.0° $\rightarrow$ output error 1.6° ($\lambda \approx 0.80$)
\end{itemize}

The contraction rate is approximately $\lambda \approx 0.8$ in the 2--5° regime, predicting convergence to a fixed point at:
\begin{equation}
    \text{ARE}^* \approx \frac{\lambda}{1-\lambda} \cdot \epsilon_{\text{model}} \approx 4 \cdot \epsilon_{\text{model}}
\end{equation}
where $\epsilon_{\text{model}}$ is Pi3X's irreducible error. With $\text{ARE}^* \approx 3.1°$, this gives $\epsilon_{\text{model}} \approx 0.8°$, consistent with the 1.2° achieved under perfect GT conditioning (which includes pipeline noise).

\paragraph{Failure for translation.} Pi3X outputs poses in Sim(3), not SE(3). The translation component $t^c$ is defined only relative to the per-run scale. Feeding translations back as SE(3) conditioning introduces a scale mismatch that does not contract:
\begin{equation}
    t^{(n+1)} = s^{(n)} \cdot f_t(R^{(n)}, t^{(n)}) + \text{bias}
\end{equation}
where $s^{(n)}$ is the unknown per-run scale. This explains why ATE degrades monotonically under na\"ive iteration.

\subsection{Denoising Score Matching Interpretation}

The Monte Carlo perturbation strategy connects to denoising score matching~\cite{vincent2011}. If Pi3X was trained with noisy pose conditioning (likely, for robustness), then conditioned inference implicitly computes:
\begin{equation}
    \mathbb{E}[R^* \mid R + \epsilon, \text{imgs}] \approx R + \sigma^2 \nabla_R \log p(R \mid \text{imgs})
\end{equation}

By running $S$ perturbations $\{R + \epsilon_s\}_{s=1}^S$ and averaging the outputs:
\begin{equation}
    \bar{R} = \text{Fr\'{e}chetMean}\left(\{f(R + \epsilon_s)\}_{s=1}^S\right)
\end{equation}
we approximate the posterior mean $\mathbb{E}[R^* \mid \text{imgs}]$ via Monte Carlo integration. This is a form of Langevin Monte Carlo on the rotation manifold, where each conditioned inference step provides a sample from the posterior.

\subsection{Connection to Diffusion Models}

The iterative refinement process resembles a single-step denoising diffusion:
\begin{enumerate}
    \item Start with a noisy estimate $R^{(0)}$ (unconditioned Pi3X output)
    \item Apply a learned denoiser $f$ conditioned on the images
    \item The denoiser moves the estimate toward the clean manifold
\end{enumerate}

Unlike standard diffusion which requires many small steps from pure noise, our setting starts close to the solution ($\sim$4° error) and requires only 2--3 denoising steps. The key difference is that our ``noise schedule'' is not designed but inherited from Pi3X's prediction uncertainty.

\section{Discussion}
\label{sec:discussion}

\paragraph{Practical deployment.} The hybrid fusion strategy (Stage~3) provides the most favourable accuracy--cost tradeoff: a single conditioned inference pass (2$\times$ compute) yields $\sim$21\% rotation improvement while preserving translation. The full MC iterative method (Stage~4, 13$\times$ compute) offers diminishing marginal gains and is appropriate primarily for offline applications where sub-3\degree{} rotation accuracy justifies the additional cost. In all cases, hybrid decoupling is essential---na\"ive iteration without it leads to progressive translation degradation regardless of the number of steps.

\paragraph{Generality of the framework.} While our experiments use Pi3X, the decoupled hybrid refinement framework applies in principle to any learned geometry model that (i) accepts pose conditioning as input, and (ii) outputs poses defined up to a gauge ambiguity. These properties are shared by several recent architectures (Table~\ref{tab:benchmark}). The contraction analysis (Section~\ref{sec:theory}) provides a diagnostic: measuring input--output error reduction under corrupted conditioning reveals whether a given model exhibits the contraction property and, if so, at what rate. Models with $\lambda < 1$ are amenable to self-refinement; those with $\lambda \geq 1$ are not.

\paragraph{Limitations.} Several constraints bound the applicability of our method. The accuracy floor ($\sim$3\degree{} ARE) is determined by the contraction rate and the model's irreducible error---it cannot be lowered without stronger external priors or architectural modifications to the conditioning mechanism. Translation accuracy is bounded by the quality of unconditioned inference, since hybrid fusion anchors translations to this source. The convergence is empirical: while iterations 1--3 consistently improve, iteration 4 and beyond can produce orientation flips (Table~\ref{tab:convergence}), necessitating a stopping criterion. Finally, our evaluation is conducted on a single multi-camera installation; generalisation to other wide-baseline deployments with different camera counts and geometries remains to be validated.

\paragraph{Future work.} Several directions merit investigation. First, Pi3X also accepts depth conditioning---monocular depth priors (e.g., Depth Anything~v2) could provide complementary geometric signal without requiring ground truth poses, potentially breaking the rotation floor. Second, partial conditioning (conditioning only high-confidence cameras while leaving others unconditioned) may exploit the model's cross-view attention to propagate reliable pose information to uncertain views. Third, the fixed perturbation magnitude ($\sigma_R = 2\degree{}$) could be replaced by an adaptive noise schedule that decreases across iterations, following the annealing principle from diffusion-based methods. Fourth, sparse geometric anchors (a pre-calibrated camera pair or LiDAR registration points) could provide a stronger initialisation that places the estimate deeper within the contraction basin, accelerating convergence and lowering the fixed point.


\bibliographystyle{plain}

\end{document}
